\Article{5}
The formula by which we expressed the product of any number of functions in the preceding article is general and assumes no relation between the roots $r,r',r'',r'''$, etc., or between $R,R',R'',R'''$, etc. We can readily deduce that, if the roots $r',r'',r'''$, etc., may be regarded as powers of $r$, and the roots $R',R'',R'''$, etc., as powers of $R$, each term of the product will have the form $R^M[r^m,R^\lambda]$, where the exponent $\lambda$ is the same for every term, namely $R^\lambda=RR'R''R'''$, etc. Thus, by what we noted in Article~3, such a product reduces to the following form:
\[
\begin{aligned}
&A[1,R^\lambda]+B[r,R^\lambda]+B'[r^g,R^\lambda]+B''[r^{g^2},R^\lambda]+B'''[r^{g^3},R^\lambda]+\text{etc.}\\
&\hspace{5em}{}+B^{(\alpha-1)}[r^{g^{\alpha-1}},R^\lambda]
\end{aligned}
\]
where each of the coefficients $A,B,B',B'',B'''$, etc., will have the form
\[
h+h'R+h''R^2+h'''R^3+\text{etc.}+h^{(\beta-1)}R^{\beta-1}
\]
and $h,h',h'',h'''$, etc., denote fixed integers.

\SourceLeaf{257}
The simplest case is that in which we put $r=r'=r''=r'''$, etc., and also $R=R'=R''=R'''$, etc.; our product then becomes the power $[r,R]^\lambda$, which therefore always reduces to the form given above.

\Article{6}
Thus, putting $\lambda=\beta$, the power $[r,R]^\beta$ takes the form
\[
\begin{aligned}
&A[1,1]+B[r,1]+B'[r^g,1]+\text{etc.}+B^{\alpha-1)}[r^{g^{\alpha-1}},1]\\
&{}=\beta\gamma A+B(\beta\gamma,1)+B'(\beta\gamma,g)+B''(\beta\gamma,g^2)+\text{etc.}\\
&\qquad{}+B^{(\alpha-1)}(\beta\gamma,g^{\alpha-1})=\theta'
\end{aligned}
\]

If, then, not only the value of the root $R$ (and hence the values of the coefficients $A,B,B'$, etc.) but also the values of the individual sums of $\beta\gamma$ terms $(\beta\gamma,1),(\beta\gamma,g)$, etc., are assumed known, the value of $\theta'$ will be known at once, and from it $[r,R]$ can be found by the formula ${}^{\beta}\!\surd\theta'$. This expression admits $\beta$ distinct values, so it might seem doubtful which one should be taken; but it is easily shown that this choice is entirely arbitrary whenever $R$ is a \emph{proper} root of the equation $x^\beta-1=0$. For in this case those $\beta$ values of the radical expression ${}^{\beta}\!\surd\theta'$ are evidently
\[
[r,R],\quad[r^{g^\alpha},R],\quad[r^{g^{2\alpha}},R]\dotsfour[r^{g^{\alpha\beta-\alpha}},R]
\]
since, by Article~3, the powers of these functions with exponent $\beta$ are equal to one another, while the functions themselves are proportional to the $\beta$ distinct roots of the equation $x^\beta-1=0$. But as long as only the sums of $\beta\gamma$ terms $(\beta\gamma,1),(\beta\gamma,g)$, etc., are known, the root $r$ itself is determined only to the extent that it must lie in the collection $(\beta\gamma,1)$, and it remains arbitrary which member of this collection we take for $r$. These roots are $r,r^{g^\alpha},r^{g^{2\alpha}}$, etc., and therefore we can also take any of the functions $[r,R],[r^{g^\alpha},R],[r^{g^{2\alpha}},R]$, etc., for $[r,R]$.

These conclusions would not hold if $R$ were not a proper root of the equation $x^\beta-1=0$. For if $R$ is assumed to be a proper root of the equation $x^{\beta'}-1=0$, with $\beta'$ a divisor of $\beta$, it is readily seen that
\[
[r,R]=[r^{g^{\alpha\beta'}},R],\quad[r^{g^\alpha},R]=[r^{g^{\alpha\beta'+\alpha}},R]\text{ etc.}
\]
so that in the collection of $\beta$ functions $[r,R],[r^{g^\alpha},R]\ldots[r^{g^{\alpha\beta-\alpha}},R]$ only $\beta'$ distinct functions occur, and therefore among the values of the expression ${}^{\beta}\!\surd\theta'$ no more than $\beta'$ are admissible, while the remaining $\beta-\beta'$ are extraneous. It is readily seen, however, that in this case we need not go as far as the power with exponent $\beta$ of the function $[r,R]$
\SourceLeaf{258}
since already the power $[r,R]^{\beta'}$ reduces to our form
\[
\beta\gamma A+B(\beta\gamma,1)+B'(\beta\gamma,g)+B''(\beta\gamma,g^2)\text{ etc.}
\]
Thus we shall obtain $[r,R]$ by an expression such as ${}^{\beta'}\!\surd\theta$, and it will make no difference which value of this expression we take.

\Article{7}
In the same way as $[r,R]$, we can also determine the functions $[r,R^2]$, $[r,R^3]$, etc., or in general $[r,R^k]$. For if substituting in $\theta'$, in place of $R$, the powers $R^2,R^3$, etc., $R^k$, is assumed to give the functions $\theta'',\theta'''$, etc., $\theta^{(k)}$, then evidently $[r,R^2]^\beta=\theta''$, $[r,R^3]^\beta=\theta'''$, etc., and in general $[r,R^k]^\beta=\theta^{(k)}$; hence these functions too can be expressed by radicals, $[r,R^2]={}^{\beta}\!\surd\theta''$, etc. But it is not suitable to use these radical expressions when a quantity is to be expressed by a function of $[r,R]$, $[r,R^2]$, etc. For since their individual values are not fully determined, it would remain doubtful which values could be combined with one another; yet this is evidently by no means arbitrary. It is readily seen that, as soon as a definite value is taken for $[r,R]$, all of $[r,R^2]$, $[r,R^3]$, etc., must also acquire fully determined values, which the radical expressions do not indicate. We must therefore set these expressions aside and seek others by which $[r,R^2]$, $[r,R^3]$, etc., can be expressed \emph{rationally} in terms of $[r,R]$ and known quantities. We readily accomplish this as follows.

By the theorem of Article~4 and what we showed in Article~5, the product $[r,R^k]\times[r,R]^{\beta-k}$ also reduces to a form such as
\[
\begin{aligned}
&\beta\gamma A+B(\beta\gamma,1)+B'(\beta\gamma,g)+B''(\beta\gamma,g^2)+\text{etc.}\\
&\qquad{}+B^{(\alpha-1)}(\beta\gamma,g^{\alpha-1})
\end{aligned}
\]
where $A,B,B',B''$, etc., are rational functions of $R$. Thus, putting the products
\[
\begin{aligned}
\relax[r,R^2]\times[r,R]^{\beta-2}&=\vartheta''\\
\relax[r,R^3]\times[r,R]^{\beta-3}&=\vartheta'''\\
\relax[r,R^4]\times[r,R]^{\beta-4}&=\beta''''\\
&\text{etc.}
\end{aligned}
\]
the quantities $\vartheta'',\vartheta''',\vartheta''''$, etc., can also be determined rationally, and
\SourceLeaf{259}
\[
\begin{aligned}
\relax[r,R^2]&=\frac{\vartheta''}{\theta'}[r,R]^2\\
\relax[r,R^3]&=\frac{\vartheta'''}{\theta'}[r,R]^3\\
\relax[r,R^4]&=\frac{\vartheta''''}{\theta'}[r,R]^4\\
&\text{etc.}
\end{aligned}
\]

These expressions therefore give the values of the functions $[r,R^2]$, $[r,R^3]$, etc., rationally, provided that we do not have $[r.R]=0$; if this equality held, they would become indeterminate. But we can rigorously prove that $[r,R]=0$ can never occur whenever $r$ denotes a root different from $1$, although we must reserve the presentation of this proof for another occasion, so as not to become too lengthy here.

\Article{8}
What we set out in the preceding articles is useful when the aim is to descend from periods of $\beta\gamma$ terms to periods of $\gamma$ terms. For it is readily seen that, if $R$ denotes a proper root, we have
\[
\begin{aligned}
\beta(\gamma,1)&=(\beta\gamma,1)+[r,R]+[r,R^2]+[r,R^3]+\text{etc.}+[r,R^{\beta-1}]\\
\beta(\gamma,g^\alpha)&=(\beta\gamma,1)+R^{\beta-1}[r,R]+R^{\beta-2}[r,R^2]\\
&\qquad{}+R^{\beta-3}[r,R^3]+\text{etc.}+R[r,R^{\beta-1}]\\
\beta(\gamma,g^{2\alpha})&=(\beta\gamma,1)+R^{2\beta-2}[r,R]+R^{2\beta-4}[r,R^2]\\
&\qquad{}+R^{2\beta-6}[r,R^3]+\text{etc.}+R^2[r,R^{\beta-1}]\\
&\text{etc.}
\end{aligned}
\]

If here, for the individual $[r,R]$, $[r,R^2]$, etc., we took the radical expressions ${}^{\beta}\!\surd\theta'$, ${}^{\beta}\!\surd\theta''$, etc., the value of each series would be uncertain among $\beta^{\beta-1}$ values; by contrast, when the rational expressions for $[r,R^2]$, etc., are taken, it will be subject to no ambiguity other than that which is unavoidable by the nature of the matter. This observation seems to have escaped the attention of the illustrious \textsc{Lagrange}, who thought he had \emph{simplified} our method given in Article~360 of the Disquis.\ arithm., where we had deliberately passed over radical expressions and proposed only rational ones, by substituting the former for the latter (Traité de la résolution numérique des équations; edition $2$, p.\ 311).

Moreover, it hardly needs mentioning here that, as soon as the values of the periods $(\gamma,1)$, $(\gamma,g^\alpha)$, etc., or of just one of them, have been found, the values of all the remaining periods of $\gamma$ terms can be deduced rationally from them. Thus the descent from periods of $\beta\gamma$ terms to periods of $\gamma$ terms requires solving the equations $x^\beta=1$, $x^\beta=\theta'$, and the remaining operations are performed rationally.
